\section{Analyse et conception}

\subsection{Partie 1 : registre à décalage à droite 4 bits}

Si \texttt{en\_sh\_i} = 1, le registre se décale à droite à chaque flanc montant de l'horloge et \texttt{ser\_r\_i} entre en $Q_3$. Sinon, il garde son état. Chaque cellule est donc une bascule D précédée d'un multiplexeur qui choisit entre maintien (entrée 0, sa propre sortie) et décalage (entrée 1, la sortie de la cellule voisine). Le reset est asynchrone et agit directement sur les bascules.

\extrait{circ_p1_cellule.png}{3cm}{Une cellule du registre 4 bits : multiplexeur et bascule D}{fig:circ_p1_cellule}

\subsection{Partie 2 : registre 8 bits avec chargement et décalage}

Le registre a trois fonctions synchrones, par ordre de priorité : chargement, décalage (à droite ou à gauche) et maintien. Le tableau~\ref{tab:fct_sync} donne l'état futur $Q^+$ du registre au prochain flanc montant de \texttt{clock\_i}, en fonction des entrées de commande. Le reset n'y figure pas car il est asynchrone : tant que \texttt{reset\_i} = 1, le registre vaut 00000000, quelles que soient l'horloge et les autres entrées.

\begin{table}[H]
  \centering
  \small
  \begin{tabular}{c c c l c}
    \hline
    \texttt{load\_i} & \texttt{en\_sh\_i} & \texttt{Right\_nLeft\_i} & \textbf{Fonction} & $Q^+_7 \dots Q^+_0$ \\
    \hline
    1 & - & - & Chargement parallèle & \texttt{val\_i(7..0)} \\
    0 & 1 & 1 & Décalage à droite    & \texttt{ser\_r\_i}, $Q_7 \dots Q_1$ \\
    0 & 1 & 0 & Décalage à gauche    & $Q_6 \dots Q_0$, \texttt{ser\_l\_i} \\
    0 & 0 & - & Maintien             & $Q_7 \dots Q_0$ \\
    \hline
  \end{tabular}
  \caption{Table des fonctions synchrones de \texttt{srg8\_ld\_sh} (\texttt{-} : valeur indifférente)}
  \label{tab:fct_sync}
\end{table}

\texttt{load\_i} n'est soumis à aucune autre condition : le chargement l'emporte donc sur le décalage. \texttt{Right\_nLeft\_i} n'est pris en compte que si \texttt{en\_sh\_i} = 1. Cette table se traduit directement par trois multiplexeurs 8 bits en cascade placés devant le registre :
\begin{itemize}
  \item le premier (\texttt{Right\_nLeft\_i}) choisit entre la valeur décalée à droite et celle décalée à gauche ;
  \item le deuxième (\texttt{en\_sh\_i}) choisit entre cette valeur décalée et l'état actuel (maintien) ;
  \item le dernier (\texttt{load\_i}) impose \texttt{val\_i}, ce qui rend le chargement prioritaire.
\end{itemize}
Le reset asynchrone passe avant tout. L'entrée WE du registre est fixée à 1.

\extrait{circ_p2_mux.png}{4cm}{Chaîne de multiplexeurs devant le registre 8 bits}{fig:circ_p2_mux}

\subsection{Partie 3 : chenillard}

Le chenillard fait tourner le pattern de manière cyclique : le bit qui sort d'un côté du registre doit rentrer de l'autre. Pour cela, $Q_0$ est rebouclé sur \texttt{ser\_r\_i} (rotation à droite) et $Q_7$ sur \texttt{ser\_l\_i} (rotation à gauche). Une porte NOT convertit \texttt{nReset\_i}, actif bas, en reset actif haut pour le registre.

\extrait{circ_p3_rebouclage.png}{5cm}{Rebouclage de $Q_0$ et $Q_7$ sur les entrées série}{fig:circ_p3_rebouclage}
